\documentclass[10pt,a4paper]{amsart}
\usepackage[margin=19mm]{geometry}
\usepackage{amsmath,amssymb,mathtools}
\usepackage[scr=rsfs,cal=euler]{mathalfa}
\usepackage{xfrac}
\usepackage[unicode,hidelinks,bookmarks=false]{hyperref}
\newcommand{\ie}{i.e.\ }
\newcommand{\cf}{cf.\ }
\newcommand{\resp}{respectively\ }
\newcommand{\Subsection}{\subsection}
\providecommand{\nobreakdash}{\nobreak\hbox{-}}
\setlength{\emergencystretch}{2em}
\multlinegap0pt
\usepackage{bm}
\usepackage{bbm}
\usepackage{dsfont}
\usepackage{xspace}
\usepackage{cancel}
\usepackage{mathtools}
\usepackage{amsthm}
\makeatletter
\@ifundefined{lem}{\newtheorem{lem}{Lem}}{}
\@ifundefined{lemma}{\newtheorem{lemma}[lem]{Lemma}}{}
\makeatother
\title[On Carmichael numbers of the form $2^np^m+1$]{On Carmichael numbers of the form $2^np^m+1$}
\author{Florian Luca}
\date{Source: arXiv:2606.28829v1 (CC BY 4.0)}
\begin{document}
\maketitle

\noindent\emph{Source and licence.} On Carmichael numbers of the form $2^np^m+1$. Version arXiv:2606.28829v1 (CC BY 4.0), distributed under the Creative Commons Attribution 4.0 International licence, as recorded in the arXiv licence metadata for this exact version and shown on the abstract page https://arxiv.org/abs/2606.28829v1. Full rights notice: licenses/arxiv-2606.28829-CC-BY-4.0.txt.\par\smallskip

\noindent\textit{Editorial scope.} Counted unit: the Lemma whose source label is \texttt{lem:12} in the article (source0.tex lines 364--368, source label \texttt{lem:12}), delivered together with the article's own complete original proof of it (source0.tex lines 371--475, one \texttt{proof} environment of 105 lines). No mathematics is written, repaired or re-derived: statement and proof are the article's own text line for line. Required context retained uncounted: eq:0 (lines 131--134); eq:102 (lines 198--201); eq:lambda (lines 227--230); lem:2 (lines 333--337). Every definition, display and label of those lines is retained unchanged. Omitted from this package and not redistributed: 1--130, 135--197, 202--226, 231--332, 338--363, 369--370, 476--731. Nothing inside a retained excerpt is omitted or altered: every retained body is delivered exactly as the article's lines are written, so no omission receipt of any kind accompanies this package. Citations: the retained excerpts carry no citation command at all, so no bibliography entry of the article is transcribed and no reference is redistributed. Reference adaptation: none was needed and none was made; every reference inside the delivered unit resolves inside it, so no unresolved reference or citation is produced. Printed slips kept literal: no printed word, symbol, hypothesis or number of the counted statement or of its proof was altered. Layout and class: the article's own class is \texttt{article} with packages including amsmath, amssymb, amsbsy, amsfonts, amsthm, latexsym, amsopn, amstext, amsxtra, euscript, amscd; the wrapper is the lane's pinned \texttt{amsart} editorial wrapper at 10pt on A4, which restates the theorem-like environments the retained text needs and adds the article's own macro definitions that the retained text needs (none), with extra wrapper packages (bm, bbm, dsfont, xspace, cancel, mathtools). The delivered numbering is the unit's own. Assets: no retained span contains a figure environment, a graphics-inclusion command, a PDF-page inclusion, a table or a TikZ picture; no image file is copied into this package, the package declares no asset dependency and no image file is redistributed with it. Licence: version arXiv:2606.28829v1 of this article is distributed under the Creative Commons Attribution 4.0 International licence, as recorded in the arXiv licence metadata for this exact version and shown on the abstract page https://arxiv.org/abs/2606.28829v1. A full rights notice is delivered at licenses/arxiv-2606.28829-CC-BY-4.0.txt. Characters: every retained excerpt is pure ASCII, so the delivered engine, XeLaTeX with its default Latin Modern text font, reports no missing character.
\begin{equation}
\label{eq:0}
N=\prod_{i=1}^k (2^{a_i}p^{b_i}+1),\qquad 1\le a_i\le n,~0\le b_i\le m,
\end{equation}

\begin{equation}
\label{eq:102}
p\mid \sum_{i\in I(b)} 2^{a_i},\qquad {\text{\rm or}}\qquad p\mid 2^n-\sum_{i\in I(b)} 2^{a_i},
\end{equation}

\begin{equation}
\label{eq:lambda}
2^{\lambda_{i}}\equiv \varepsilon_{i}\pmod {q_i},
\end{equation}

\begin{lemma}
\label{lem:2}
Assume that $N=2^np^m+1$ is Carmichael for some odd $m\ge 5$ and prime $p>500m^5\log m$. Let $\gamma:=\min\{a_i: q_i\mid N\}$. 
Then $\gamma<n$. Further, there are an even number of prime factors $q_i$ of $N$ with $a_i=\gamma$. Additionally, $b\le m-1$, so divisibility \eqref{eq:102}-right cannot occur. Furthermore,  $\gamma\ge 5$. 
\end{lemma}

\begin{lemma}
\label{lem:12}
If $N=2^np^m+1$ is Carmichael, $m\ge 5$ is odd, $p>500m^5\log m$ is prime, then $n=2^{n_0}n_1$, where $n_1$ is odd and $2^{n_0}>(\log p)/m^3$. Additionally,  for every prime factor $q_i=2^{a_i}p^{b_i}+1$ of $N$, we have $a_i=2^{\beta_i}\lambda_i$ where $\lambda_i$ is odd and 
$2^{\beta_i}>(\log p)/m^3$. 
\end{lemma}

\begin{proof}
Assume $q_i=2^{a_i}p^{b_i}+1$ is a prime factor of $N$ with $b_i>0$. Let ${\text{\rm ord}}_2(q_i)=2^{\omega_i}$. We then have 
$$
2^{2^{\omega_i-1}}\equiv -1\pmod {q_i}.
$$
Thus, 
\begin{equation}
\label{e:10}
2^{2^{\omega_i-1}}+1\ge 2^{a_i}p^{b_i}+1>p+1,\quad {\text{\rm therefore}}\quad 2^{\omega_i-1}\ge \frac{\log p}{\log 2}>\log p.
\end{equation}
We shall exploit the above inequality for various prime factors $q_i$ of $N$ with $b_i>0$. We write $a_i=2^{\beta_i}\gamma_i$ for some  nonzero negative integers $\beta_i$ and odd integers $\gamma_i$. 
Returning to congruences \eqref{eq:lambda}, we get that 
$$
2^{\lambda_i}\equiv \varepsilon_i\pmod {q_i},
$$
where $\lambda_i=nc_i-a_im_i$ and $\varepsilon_i\in \{\pm 1\}$. It follows that 
\begin{equation}
\label{eq:100}
nc_i-a_im_i=2^{\omega_i-1}\delta_i,\quad {\text{\rm holds~with~some~nonzero~integer}}\quad \delta_i.
\end{equation}
The fact that $\delta_i$ is nonzero follows from the fact that $2^np^m$ and $2^{a_i}p^{b_i}$ are multiplicatively independent for all prime factors $q_i=2^{a_i}p^{b_i}+1$ with $b_i>0$ of $N$. 
We now calculate $2$-adic valuations in \eqref{eq:100} above. Write $n=2^{n_0}n_1$ with $n_1$ odd. Assume first that 
\begin{equation}
\label{eq:11}
2^{n_0}\ge\frac{\log p}{m^3}. 
\end{equation}
Note that since $m_i\mid m$ is odd, it follows that $\nu_2(a_im_i)=\nu_2(a_i)=\beta_i$. Thus, 
we get that 
$$
\beta_i=\nu_2(a_im_i)\ge \min\{\nu_2(nc_i),\nu_2(2^{\omega_i-1}\delta_i)\}\ge \min\{n_0,\omega_i-1\},
$$
therefore 
$$
2^{\beta_i}\ge \min\{2^{n_0},2^{\omega_i-1}\}\ge \frac{\log p}{m^3}.
$$
So, from now on we shall assume that \eqref{eq:11} fails. We take the $t\ge 2$ primes of the form $q_i=2^{\gamma}p^{b_i}+1$ which are guaranteed by Lemma \ref{lem:2}. Assume that there are two of them say $q_i=2^{\gamma}p^{b_i}+1$ and $q_j=2^{\gamma}p^{b_j}+1$ with $b_ib_j> 0$. We then get that 
$$
nc_i-\gamma m_i=2^{\omega_i-1}\delta_i\qquad {\text{\rm and}}\qquad nc_j-\gamma m_j=2^{\omega_j-1}\delta_j.
$$
Eliminating $\gamma$ from the above two equations we have 
$$
n(m_jc_i-m_ic_j)=2^{\min\{\omega_i-1,\omega_j-1\}}\Lambda
$$
for some integer $\Lambda$. Note that $m_ic_j\ne m_jc_i$. Indeed, if they are equal then $m_i/c_i=m_j/c_j$. But $m_i/c_i=m/b_i$ and $m_j/c_j=m/b_j$. So, these fractions being equal entails $b_i=b_j$, so $q_i=q_j$ 
which is false. Now 
$$
|m_ic_j-m_jc_i|<m^2,
$$
which shows that $\nu_2(|m_ic_j-m_jc_i|)<2\log m/\log 2$. This shows that 
$$
2^{n_0}> \frac{2^{\min\{\omega_i-1,\omega_j-1\}}}{m^2}>\frac{\log p}{m^2},
$$
contradicting \eqref{eq:11}. Thus, we have to assume that $t=2$, and up to relabelling the primes we have $q_1=2^{\gamma}+1$, $\gamma=2^{\beta}$, and $q_2=2^{\gamma}p^{b_2}+1$ for some $b_2>0$. We write again 
\begin{equation}
\label{eq:12}
nc_2-\gamma m_2=2^{\omega_2-1}\delta_2.
\end{equation}
Assume next that $b_2$ is odd. Then $c_2\mid b_2$ is also odd, so 
$$
\nu_2(nc_2)=\nu_2(n)=\nu_2(\gamma m_2)=\nu_2(\gamma)=\beta.
$$ 
In particular, $n_0=\beta$. Now 
$$
N=2^np^m+1=2^{\gamma n_1} p^m+1\equiv (-1)^{n_1} p^m+1\pmod {q_1}\equiv p^m-1\pmod {q_1}.
$$
Thus, $p$ has order dividing $m$ modulo $q_1$. Since $q_1$ is a Fermat prime, the order of $p$ modulo $q_1$ must be a power of $2$. Since this order divides $m$ which is odd, it must be $1$. This shows that $p\equiv 1\pmod {q_1}$. Now 
$$
q_2=2^{\gamma}p^{b_2}+1\equiv (-1) (1)^{b_2}+1\equiv 0\pmod {q_1},
$$
which shows that $q_1\mid q_2$, which is of course impossible since $q_1$ and $q_2$ are distinct primes. Thus, $b_2$ is even. Let $a_2:=\min\{a_{s}: s\ge 3\}$. Note that $a_2$ exists since 
$k=\omega(N)\ge 3$. Reducing relation \eqref{eq:0} modulo $2^{a_2}$ we get that
$$
1\equiv (2^{\gamma}p^{b_2}+1)(2^{\gamma}+1)\equiv 1+2^{\gamma}(p^{b_2}+1)+2^{2\gamma}p^{b_2}\pmod {2^{a_2}},
$$
which shows that $2^{\min\{a_2,2\gamma\}}\mid 2^{\gamma}(p^{b_2}+1)$. Since $b_2$ is even, we get that $\min\{a_2,2\gamma\}=\gamma+1$. Since $\gamma\ge 8$, we get that $a_2=\gamma+1$. In particular, 
$a_2$ is odd. Let $q_3=2^{\gamma+1}p^{b_3}+1$. Clearly, $b_3\ne 0$ since otherwise $q_3$ should be a Fermat prime which is false since $\gamma+1\ge 9$ is odd. Thus, $b_3\ge 1$. Note also that $b_3$ is odd since
if $b_3$ is even then $q_3=2^{\gamma+1}(p^2)^{b_3/2}+1\equiv 0\pmod 3$, so $q_3$ is not prime. We now write 
\begin{equation}
\label{eq:13}
nc_3-(\gamma+1)m_3=2^{\omega_3-1}\delta_3.
\end{equation}
Note that $(\gamma+1)m_3$ is odd showing that $nc_3$ is also odd. In particular, $n$ is odd. Since by \eqref{eq:12} we have that $\nu_2(nc_2)=\nu_2(\gamma m_2)=\nu_2(\gamma)=\beta$, we get that 
$\beta=\nu_2(\gamma)=\nu_2(c_2)\le \log m/\log 2$. Thus, $\gamma=2^{\beta}\le m-1$. We now eliminate the terms containing $\gamma$ in \eqref{eq:12} and \eqref{eq:13} getting 
\begin{equation}
\label{eq:14}
n(c_2 m_3(\gamma+1)-c_3m_2\gamma)=2^{\min\{\omega_2-1,\omega_3-1\}}\Lambda',
\end{equation}
for some integer $\Lambda'$. Let us prove that the factor multiplying $n$ above is nonzero. If it were zero, we would get
$$
\frac{c_2(\gamma+1)}{m_2}=\frac{c_3\gamma}{m_3}.
$$
Since $c_2/m_2=b_2/m$ and $c_3/m_3=b_3/m$, we get that $(\gamma+1)b_2=\gamma b_3$, so $b_2=\gamma d$ and $b_3=(\gamma+1)d$ for some integer $d$. Thus, 
$$
q_2=2^{\gamma} p^{b_2}+1=(2p^d)^{\gamma}+1,\qquad q_3=2^{\gamma+1} p^{b_3}+1=(2p^d)^{\gamma+1}+1,
$$
and we see that $q_3$ is divisible by $2p^d+1$ (since $\gamma+1>1$ is odd), so in particular it is not prime. Thus, in \eqref{eq:13}, the factor multiplying $n$ is not zero. Its size is at most
$$
|c_2m_3(\gamma+1)-c_3m_2\gamma|\le \max\{c_2,c_3\}\max\{m_2,m_3\}(\gamma+1)\le m^3,
$$
which shows that 
$$
2^{n_0}> \frac{2^{\min\{\gamma_2-1,\gamma_3-1\}}}{m^3}\ge \frac{\log p}{m^3},
$$
so in fact \eqref{eq:11} holds. The lemma is therefore proved.  
\end{proof}

\end{document}
